\SourceLeaf{24}
\[
[6]\qquad X=\frac{4a\pi BC}{AA}\int\frac{tt\mathrm{d}t}{\BareRoot{((1-(1-\frac{BB}{AA})tt)(1-(1-\frac{CC}{AA})tt)}}
\]
wobei von $t=0$ bis $t=1$ integriert wird. Die zu den Koordinatenachsen $y,z$ parallelen Anziehungen ergeben sich daraus offenbar unmittelbar, wenn $a,A$ mit $b,B$ oder mit $c,C$ vertauscht wird.

Diese Formel liefert also die Anziehung auf alle Punkte im Inneren des Sphäroids. Da sie streng gilt, wie nahe der angezogene Punkt der Oberfläche des Sphäroids auch kommt, gilt sie auch für Punkte auf der Oberfläche. Die Anziehung auf äußere Punkte ist bereits auf die Anziehung auf Punkte dieser Oberfläche zurückgeführt worden, sodass die Aufgabe nun vollständig gelöst ist.

Gleichung $[6]$ zeigt ferner, dass für einen inneren Punkt die Anziehung aller ähnlichen und ähnlich gelegenen Sphäroide genau gleich ist. Man denke sich daher ein solches Sphäroid in mehrere Schichten geteilt, deren innere und äußere Flächen der Oberfläche des Sphäroids ähnlich sind und ähnlich liegen. Die einzelnen Schichten, die den angezogenen Punkt umgeben, tragen offenbar nichts zur Anziehung auf diesen Punkt bei, sodass nur die Anziehung des inneren Kerns übrig bleibt. Seine Oberfläche geht durch den Punkt selbst.

\Article{14}
\par
Die Integration der Formel $[6]$ selbst braucht nicht ausführlich erörtert zu werden. Bekanntlich hängt sie von höheren transzendenten Funktionen als Kreisfunktionen und Logarithmen ab, wenn $A,B,C$ alle ungleich sind. In diesem Fall greifen wir daher zu Reihen, die umso schneller konvergieren, je weniger das Sphäroid von der Kugel abweicht. Sind aber zwei der Größen $A,B,C$ gleich, beispielsweise $A=B$, sodass das Sphäroid durch Drehung um die Achse $=2C$ entsteht, so gilt
\[
\begin{aligned}
X&=\frac{4\pi aC}{A}\int\frac{tt\mathrm{d}t}{\BareRoot{(1-(1-\frac{CC}{AA})tt)}}\\
&=\frac{2\pi a\cos\varphi}{\sin\varphi^3}(\varphi-\tfrac12\sin2\varphi)
\end{aligned}
\]
wenn man $\frac{C}{A}=\cos\varphi$ oder $\BareRoot{(1-\frac{CC}{AA})}=\sin\varphi$ setzt, falls $C<A$, oder
\[
X=\frac{2\pi aCC}{CC-AA}-\frac{2\pi aAAC}{(CC-AA)^{\frac32}}\log\frac{C+\BareRoot{(CC-AA)}}{A}
\]
falls $C>A$.
